\documentclass[12pt,a4paper]{article}
\usepackage[utf8]{inputenc}
\usepackage[T2A,T1]{fontenc}
\usepackage[ukrainian]{babel}
\usepackage{amsmath,amssymb}
\usepackage{booktabs}
\usepackage{longtable}
\usepackage{array}
\usepackage{calc}
\usepackage{geometry}
\geometry{margin=2.5cm}

\begin{document}
\section{ФОРМУЛА ВИНАХОДУ}\label{ux444ux43eux440ux43cux443ux43bux430-ux432ux438ux43dux430ux445ux43eux434ux443}



\subsection{1. Система обробки EEG-даних та обчислення CME}\label{ux441ux438ux441ux442ux435ux43cux430-ux43eux431ux440ux43eux431ux43aux438-eeg-ux434ux430ux43dux438ux445-ux442ux430-ux43eux431ux447ux438ux441ux43bux435ux43dux43dux44f-cme}



\textbf{1.} Система обробки багатоканальних EEG-даних для обчислення показника обчислювальної ментальної енергії (CME), що містить:



\begin{enumerate}

\def\labelenumi{(\roman{enumi})}

\item

  пристрій реєстрації нейросигналів EEG (110), виконаний з можливістю реєстрації біоелектричної активності мозку щонайменше по чотирьох електродах;

\item

  пристрій стримінгу (115), з'єднаний з пристроєм реєстрації (110) через Bluetooth та з модулем приймання даних (120) через мережевий протокол, виконаний з можливістю приймання сирих даних EEG, буферизації та передачі на сервер;

\item

  модуль приймання даних (120), з'єднаний з пристроєм стримінгу (115) та виконаний з можливістю зчитування потоків багатоканальних EEG-даних із часовими мітками та метаданими;

\item

  модуль сегментації (130), з'єднаний з модулем приймання даних (120) та виконаний з можливістю поділу безперервного потоку EEG-даних на дискретні часові вікна фіксованої тривалості \(\Delta\);

\item

  модуль попередньої обробки (140), з'єднаний з модулем сегментації (130) та виконаний з можливістю фільтрації, нормалізації сигналу, пригнічення мережевих завад, видалення дрейфу, контролю якості контакту електродів та маркування артефактів;

\item

  модуль виділення спектральних ознак (150), з'єднаний з модулем попередньої обробки (140) та виконаний з можливістю обчислення спектральної потужності \(P_{\text{band}}(ch, t)\) у стандартних частотних діапазонах \(\delta\), \(\theta\), \(\alpha\), \(\beta\), \(\gamma\) для кожного електроду \(ch\) у кожному вікні \(t\);

\item

  модуль формування вектора ознак (160), з'єднаний з модулем виділення спектральних ознак (150) та виконаний з можливістю формування вхідного вектора \(\mathbf{x}_t\), що містить спектральні потужності по кожному електроду та діапазону, індекс складності задачі \(c(t)\) та індикатори якості сигналу \(q(t)\);

\item

  квантовий обчислювальний модуль (170), з'єднаний з модулем формування вектора ознак (160) та виконаний з можливістю зведення повного вектора \(\mathbf{x}_t\) до квантового підвектора \(\mathbf{x}_t^{(q)} \in \mathbb{R}^8\) (усереднені потужності по діапазонах, фронтальна \(\alpha\)-асиметрія, індекс залученості \(\beta/\theta\), складність \(c(t)\)), кодування \(\mathbf{x}_t^{(q)}\) у квантовий стан за допомогою варіаційної квантової схеми з параметрами \(\Theta\) та глибиною \(D\), виконання цієї схеми на квантовому процесорі або симуляторі з кількістю вимірювань \(S\) та отримання оцінки ймовірності цільового стану \(p_{\text{flow}}(t)\);

\item

  модуль обчислення CME (180), з'єднаний з модулем виділення спектральних ознак (150), модулем формування вектора ознак (160) та квантовим обчислювальним модулем (170) та виконаний з можливістю обчислення миттєвої інтенсивності та показника обчислювальної ментальної енергії у вікні \(t\) за формулою:

\end{enumerate}



\[\text{CME}_{\text{rate}}(t) = \kappa \cdot E_{\text{band}}(t) \cdot g\bigl(c(t),\;p_{\text{flow}}(t)\bigr), \quad \text{CME}(t) = \text{CME}_{\text{rate}}(t) \cdot \Delta\]



де \(\kappa\) --- безрозмірний коефіцієнт масштабування (визначається з калібрувального набору або персонального профілю), \(E_{\text{band}}(t)\) --- зведена спектральна потужність (мкВ²), \(g(c, p)\) --- безрозмірна функція інтеграції складності та стану потоку (flow), \(\Delta\) --- тривалість вікна (с); одиниці \(\text{CME}(t)\): Вн (\(1\;\text{Вн} \equiv 1\;\text{мкВ}^2 \cdot \text{с}\));



\begin{enumerate}

\def\labelenumi{(\alph{enumi})}

\setcounter{enumi}{23}

\tightlist

\item

  модуль метаевристичної оптимізації (190), з'єднаний з квантовим обчислювальним модулем (170) та модулем обчислення CME (180) та виконаний з можливістю підбору параметрів \(\Theta\) квантової схеми та ресурсних параметрів \(S\), \(D\) за цільовою функцією, що враховує якість оцінювання та обчислювальну вартість;

\end{enumerate}



\begin{enumerate}

\def\labelenumi{(\roman{enumi})}

\setcounter{enumi}{10}

\item

  сховище (200), з'єднане з модулем обчислення CME (180) та виконане з можливістю зберігання результатів обчислення CME, журналів запитів, параметрів моделі та метрик сесій;

\item

  інтерфейс інтеграції (210), з'єднаний зі сховищем (200) та виконаний з можливістю надання доступу до результатів обчислення через програмний інтерфейс та/або веб-інтерфейс,

\end{enumerate}



яка відрізняється тим, що квантовий обчислювальний модуль (170) виконаний у вигляді 4-кубітної варіаційної квантової схеми з архітектурою повторного завантаження даних (data re-uploading), що містить \(L + 1\) повторюваних шарів, кожен з яких включає: двоканальне кутове кодування 8 ознак через обертання \(R_y\) та \(R_z\), шар ZZ-взаємодії ознак для кодування попарних кореляцій, кільцеве заплутування CNOT-гейтами та варіаційний анзац із обертаннями \(R_y\), \(R_z\) з навченими параметрами \(\Theta \in \mathbb{R}^{8(L+1)}\), а модуль метаевристичної оптимізації (190) виконує ітераційний підбір параметрів \(\Theta\), кількості shots \(S\) та глибини \(D\) квантової схеми за цільовою функцією:



\[J(\mathbf{u}) = \mathcal{L}_{\text{val}}(\mathbf{u}) + \xi_S \cdot S + \xi_D \cdot D + \xi_T \cdot T_{\text{qpu}}(\mathbf{u})\]



де \(\mathcal{L}_{\text{val}}(\mathbf{u})\) --- помилка оцінювання на валідаційній вибірці, \(T_{\text{qpu}}(\mathbf{u})\) --- час виконання на QPU, \(\xi_S, \xi_D, \xi_T \geq 0\) --- вагові коефіцієнти штрафу за ресурси, із використанням щонайменше одного з метаевристичних алгоритмів: генетичного алгоритму, оптимізації роєм частинок, мурашиного алгоритму або імітації відпалу.



\begin{center}\rule{0.5\linewidth}{0.5pt}\end{center}



\subsection{2. Спосіб обробки EEG-даних та обчислення CME}\label{ux441ux43fux43eux441ux456ux431-ux43eux431ux440ux43eux431ux43aux438-eeg-ux434ux430ux43dux438ux445-ux442ux430-ux43eux431ux447ux438ux441ux43bux435ux43dux43dux44f-cme}



\textbf{2.} Спосіб обробки багатоканальних EEG-даних для обчислення показника обчислювальної ментальної енергії (CME), що включає кроки:



\begin{enumerate}

\def\labelenumi{(\roman{enumi})}

\item

  отримання сирих даних EEG від пристрою реєстрації через пристрій стримінгу та приймання потоку багатоканальних EEG-даних від щонайменше чотирьох електродів із часовими мітками та метаданими;

\item

  сегментацію прийнятого потоку EEG-даних на дискретні часові вікна фіксованої тривалості \(\Delta\);

\item

  попередню обробку сегментованих даних, що включає фільтрацію, нормалізацію, пригнічення мережевих завад, видалення дрейфу та маркування артефактів;

\item

  обчислення спектральної потужності \(P_{\text{band}}(ch, t)\) у частотних діапазонах \(\delta\) (1--4 Гц), \(\theta\) (4--8 Гц), \(\alpha\) (8--13 Гц), \(\beta\) (13--30 Гц) та опційно \(\gamma\) (30--45 Гц, за наявності апаратної підтримки) для кожного електроду \(ch\) у кожному вікні \(t\);

\item

  обчислення зведеної спектральної потужності \(E_{\text{band}}(t)\) (мкВ²) як зваженої суми спектральних потужностей:

\end{enumerate}



\[E_{\text{band}}(t) = \sum_{ch \in \mathcal{CH}} \Bigl[\,w_\delta \cdot P_\delta(ch,t) + w_\theta \cdot P_\theta(ch,t) + w_\alpha \cdot P_\alpha(ch,t) + w_\beta \cdot P_\beta(ch,t) + w_\gamma \cdot P_\gamma(ch,t)\,\Bigr]\]



де \(w_\delta\), \(w_\theta\), \(w_\alpha\), \(w_\beta\), \(w_\gamma\) --- вагові коефіцієнти;



\begin{enumerate}

\def\labelenumi{(\roman{enumi})}

\setcounter{enumi}{5}

\item

  формування вхідного вектора ознак \(\mathbf{x}_t\), що містить спектральні потужності по електродах та діапазонах, індекс складності задачі \(c(t)\) та індикатори якості сигналу;

\item

  зведення вектора \(\mathbf{x}_t\) до квантового підвектора \(\mathbf{x}_t^{(q)} \in \mathbb{R}^8\) (усереднені потужності по діапазонах, фронтальна \(\alpha\)-асиметрія, індекс залученості \(\beta/\theta\), складність \(c(t)\)) та його кодування у квантовий стан варіаційної квантової схеми з параметрами \(\Theta\);

\item

  виконання квантової схеми на квантовому процесорі або симуляторі з кількістю вимірювань \(S\) та отримання оцінки ймовірності цільового стану \(p_{\text{flow}}(t)\);

\item

  обчислення миттєвої інтенсивності та показника CME у вікні \(t\) за формулою:

\end{enumerate}



\[\text{CME}_{\text{rate}}(t) = \kappa \cdot E_{\text{band}}(t) \cdot g\bigl(c(t),\;p_{\text{flow}}(t)\bigr), \quad \text{CME}(t) = \text{CME}_{\text{rate}}(t) \cdot \Delta\]



де \(\kappa\) --- безрозмірний коефіцієнт масштабування (визначається з калібрувального набору або персонального профілю), \(g(c, p)\) --- безрозмірна функція інтеграції складності та стану потоку (flow); одиниці \(\text{CME}(t)\): Вн (\(1\;\text{Вн} \equiv 1\;\text{мкВ}^2 \cdot \text{с}\), визначення --- у описі винаходу, п. 2);



\begin{enumerate}

\def\labelenumi{(\alph{enumi})}

\setcounter{enumi}{23}

\tightlist

\item

  агрегацію показника по сесії:

\end{enumerate}



\[\text{CME}_{\text{session}} = \sum_{t} \text{CME}(t)\]



який відрізняється тим, що на кроці (vii) кодування здійснюється у \(L + 1\) повторюваних шарах, де у кожному шарі: 8 ознак кодуються через обертання \(R_y\) з кутами \((x_i + 1) \cdot \pi\) та \(R_z\) з кутами \((x_{i+4} + 1) \cdot \pi\) для кожного з 4 кубітів (двоканальне кодування), потім застосовуються ZZ-взаємодії ознак для кодування попарних кореляцій, кільцеве заплутування CNOT-гейтами та варіаційний анзац із обертаннями \(R_y(\theta_i^{(l)})\) та \(R_z(\phi_i^{(l)})\) з навченими параметрами \(\Theta \in \mathbb{R}^{8(L+1)}\), а ймовірність \(p_{\text{flow}}(t)\) на кроці (viii) обчислюється як:



\[p_{\text{flow}}(t) = \Pr\bigl[b_0 = 1\bigr] \approx \hat{p}_{\text{flow}}(t) = \frac{1}{S}\sum_{s=1}^{S}\mathbf{1}\bigl[b_0^{(s)} = 1\bigr]\]



де \(b_0^{(s)}\) --- результат вимірювання міткового кубіта \(q_0\) у \(s\)-му прогоні, що кодує клас ``flow'' (\(= 1\)) / ``no-flow'' (\(= 0\)); причому параметри \(\Theta\) оптимізуються метаевристичним алгоритмом на розмічених вікнах з міткою \(\text{FlowLabel}(t) \in \{0,\,1\}\) за цільовою функцією \(J(\mathbf{u}) = \mathcal{L}_{\text{val}}(\mathbf{u}) + \xi_S \cdot S + \xi_D \cdot D + \xi_T \cdot T_{\text{qpu}}(\mathbf{u})\), що одночасно мінімізує помилку оцінювання та обчислювальну вартість.



\begin{center}\rule{0.5\linewidth}{0.5pt}\end{center}



\subsection{3. Деталізація квантової схеми}\label{ux434ux435ux442ux430ux43bux456ux437ux430ux446ux456ux44f-ux43aux432ux430ux43dux442ux43eux432ux43eux457-ux441ux445ux435ux43cux438}



\textbf{3.} Система за п. 1, яка відрізняється тим, що варіаційна квантова схема квантового обчислювального модуля (170) складається з \(L + 1\) повторюваних шарів (за замовчуванням \(L = 2\)), де кожен шар \(l\) містить: (i) підшар двоканального кутового кодування, в якому для кожного з 4 кубітів \(i\) застосовуються гейти \(R_y\bigl((x_i + 1) \cdot \pi\bigr)\) та \(R_z\bigl((x_{i+4} + 1) \cdot \pi\bigr)\), де \(x_0, \ldots, x_7\) --- 8 нормалізованих ознак EEG з діапазону \([-1,\,1]\);



\begin{enumerate}

\def\labelenumi{(\roman{enumi})}

\setcounter{enumi}{1}

\item

  підшар ZZ-взаємодії ознак, в якому для кожної кільцевої пари \((q_i, q_{(i+1) \bmod 4})\) застосовується \(RZZ\bigl((x_i - x_{(i+1) \bmod 4})^2 \cdot \pi\bigr)\) через CNOT-RZ-CNOT декомпозицію;

\item

  підшар кільцевого заплутування із \(\text{CNOT}(q_i,\,q_{(i+1) \bmod 4})\) для \(i = 0, 1, 2, 3\);

\item

  варіаційний підшар, в якому для кожного кубіта \(i\) застосовуються гейти \(R_y(\theta_i^{(l)})\) та \(R_z(\phi_i^{(l)})\), де \(\theta_i^{(l)} = \Theta[8l + 2i]\), \(\phi_i^{(l)} = \Theta[8l + 2i + 1]\);

\item

  вимірювання всіх 4 кубітів у обчислювальному базисі після останнього шару; ймовірність стану потоку визначається як маргінальна ймовірність міткового кубіта \(q_0\): \(p_{\text{flow}}(t) \approx \frac{1}{S}\sum_{s=1}^{S}\mathbf{1}[b_0^{(s)}=1]\).

\end{enumerate}



\subsection{4. Метаевристична оптимізація}\label{ux43cux435ux442ux430ux435ux432ux440ux438ux441ux442ux438ux447ux43dux430-ux43eux43fux442ux438ux43cux456ux437ux430ux446ux456ux44f}



\textbf{4.} Система за п. 1, яка відрізняється тим, що модуль метаевристичної оптимізації (190) виконує ітераційний цикл, що включає: (i) ініціалізацію популяції з \(N\) кандидатів із випадковими параметрами \(\Theta \in [0,\,2\pi]^{8(L+1)}\), де \(L\) --- кількість шарів повторного завантаження;



\begin{enumerate}

\def\labelenumi{(\roman{enumi})}

\setcounter{enumi}{1}

\item

  для кожного покоління --- оцінку кожного кандидата шляхом виклику квантового обчислювального модуля (170) із параметрами кандидата та обчислення фітнес-значення за цільовою функцією \(J(\mathbf{u})\) (п. 1);

\item

  еволюцію популяції шляхом відбору, кросоверу та мутації, де ймовірність мутації залежить від обраного алгоритму;

\item

  збереження параметрів найкращого кандидата як активної моделі для використання при подальшому інференсі.

\end{enumerate}



\subsection{5. Функція модуляції складності та стану потоку}\label{ux444ux443ux43dux43aux446ux456ux44f-ux43cux43eux434ux443ux43bux44fux446ux456ux457-ux441ux43aux43bux430ux434ux43dux43eux441ux442ux456-ux442ux430-ux441ux442ux430ux43dux443-ux43fux43eux442ux43eux43aux443}



\textbf{5.} Спосіб за п. 2, який відрізняється тим, що функція інтеграції складності та стану потоку (flow) \(g(c, p)\) є конфігурованою безрозмірною монотонною по кожному аргументу функцією, зокрема може мати лінійно-білінійний вигляд:



\[g(c,\,p) = \lambda_1 \cdot c + \lambda_2 \cdot p + \lambda_3 \cdot c \cdot p\]



де \(\lambda_1\), \(\lambda_2\), \(\lambda_3 \geq 0\) --- конфігуровані вагові коефіцієнти, \(c \in [0,\,1]\) --- нормований індекс складності задачі, \(p \in [0,\,1]\) --- ймовірність цільового стану, оцінена квантовим модулем; значення функції \(g(c,\,p) \in [0,\,\lambda_1 + \lambda_2 + \lambda_3]\) та за потреби обмежується до \([0,\,g_{\max}]\) для стабілізації масштабу.



\subsection{6. Детектування стану потоку та метрики сесії}\label{ux434ux435ux442ux435ux43aux442ux443ux432ux430ux43dux43dux44f-ux441ux442ux430ux43dux443-ux43fux43eux442ux43eux43aux443-ux442ux430-ux43cux435ux442ux440ux438ux43aux438-ux441ux435ux441ux456ux457}



\textbf{6.} Спосіб за п. 2, який відрізняється тим, що додатково включає крок детектування стану flow у вікні \(t\) за умовою:



\[p_{\text{flow}}(t) \;\geq\; \text{FlowThreshold}\]



де \(\text{FlowThreshold}\) --- конфігурований пороговий коефіцієнт, та обчислення метрик сесії, що включають частку вікон у стані flow, тривалість найдовшої безперервної серії flow, середнє та максимальне значення CME.



\subsection{7. Варіанти архітектури розгортання}\label{ux432ux430ux440ux456ux430ux43dux442ux438-ux430ux440ux445ux456ux442ux435ux43aux442ux443ux440ux438-ux440ux43eux437ux433ux43eux440ux442ux430ux43dux43dux44f}



\textbf{7.} Система за п. 1, яка відрізняється тим, що модулі системи реалізовані з можливістю горизонтального масштабування та підтримують: (i) синхронний режим виконання в одному обчислювальному процесі із латентністю, визначеною сумою часів обробки окремих модулів;



\begin{enumerate}

\def\labelenumi{(\roman{enumi})}

\setcounter{enumi}{1}

\item

  розподілений режим виконання, у якому модулі попередньої обробки, квантового обчислення та зберігання функціонують на окремих обчислювальних вузлах, з'єднаних мережевим протоколом, із можливістю незалежного масштабування пропускної здатності кожного вузла;

\item

  асинхронний режим з буферизацією запитів у черзі повідомлень, що забезпечує згладжування пікового навантаження та відмовостійкість при недоступності квантового процесора,

\end{enumerate}



причому математична суть обчислення CME є однаковою для всіх режимів, а вибір режиму визначається вимогами до латентності, пропускної здатності та відмовостійкості.



\subsection{8. Потоковий режим реального часу}\label{ux43fux43eux442ux43eux43aux43eux432ux438ux439-ux440ux435ux436ux438ux43c-ux440ux435ux430ux43bux44cux43dux43eux433ux43e-ux447ux430ux441ux443}



\textbf{8.} Спосіб за п. 2, який відрізняється тим, що кроки (i)--(x) виконуються у потоковому режимі реального часу, де кожне нове вікно тривалістю \(\Delta\) обробляється по мірі його формування, а результат \(\text{CME}(t)\) повертається клієнту із загальною латентністю, що включає час попередньої обробки, час звернення до квантового модуля та час запису у сховище.



\subsection{9. Калібрування коефіцієнта масштабування κ}\label{ux43aux430ux43bux456ux431ux440ux443ux432ux430ux43dux43dux44f-ux43aux43eux435ux444ux456ux446ux456ux454ux43dux442ux430-ux43cux430ux441ux448ux442ux430ux431ux443ux432ux430ux43dux43dux44f-ux3ba}



\textbf{9.} Система за п. 1, яка відрізняється тим, що коефіцієнт масштабування \(\kappa\) визначається одним із способів у порядку пріоритету: (i) з персонального профілю користувача (\(\kappa_{\text{user}}\)), обчисленого за попередніми сесіями;



\begin{enumerate}

\def\labelenumi{(\roman{enumi})}

\setcounter{enumi}{1}

\tightlist

\item

  з калібрувального періоду (перші \(N_{\text{cal}}\) вікон поточної сесії) як:

\end{enumerate}



\[\kappa = \frac{C_{\text{target}}}{\displaystyle\max_{t \in T_{\text{cal}}}\;\text{CME}_{\text{rate,raw}}(t)}\]



де \(C_{\text{target}}\) --- конфігурований цільовий максимум шкали (наприклад, \(C_{\text{target}} = 100\) для шкали \([0,\,100]\)), \(\text{CME}_{\text{rate,raw}}(t) = E_{\text{band}}(t) \cdot g(c(t),\,p_{\text{flow}}(t))\) --- некаліброване значення;



\begin{enumerate}

\def\labelenumi{(\roman{enumi})}

\setcounter{enumi}{2}

\tightlist

\item

  як фіксована емпірична константа (\(\kappa_{\text{default}}\)) --- якщо персональний профіль та калібрувальні дані відсутні,

\end{enumerate}



причому в усіх випадках \(\kappa\) забезпечує міжсесійну порівнянність показника CME (Вн) та безрозмірної шкали відображення \(\text{CME}_{\text{index}} \in [0,\,100]\).



\subsection{10. Нормалізація ознак для квантової схеми}\label{ux43dux43eux440ux43cux430ux43bux456ux437ux430ux446ux456ux44f-ux43eux437ux43dux430ux43a-ux434ux43bux44f-ux43aux432ux430ux43dux442ux43eux432ux43eux457-ux441ux445ux435ux43cux438}



\textbf{10.} Спосіб за п. 2, який відрізняється тим, що на кроці (iii) нормалізація здійснюється шляхом масштабування кожної ознаки вхідного вектора до діапазону \([-1,\,1]\) за допомогою мін-макс перетворення з калібрувальними статистиками:



\[f_{\text{norm}}[i] = 2 \cdot \frac{f[i] - f_{\min}^{(i)}}{f_{\max}^{(i)} - f_{\min}^{(i)}} - 1\]



де \(f_{\min}^{(i)}\), \(f_{\max}^{(i)}\) визначаються під час калібрувального періоду або з персонального профілю, що забезпечує кути обертання у \([0,\,2\pi]\) при кодуванні \((f_{\text{norm}} + 1) \cdot \pi\) у квантовій схемі та повне покриття сфери Блоха.



\subsection{11. Класичний нейромережевий модуль}\label{ux43aux43bux430ux441ux438ux447ux43dux438ux439-ux43dux435ux439ux440ux43eux43cux435ux440ux435ux436ux435ux432ux438ux439-ux43cux43eux434ux443ux43bux44c}



\textbf{11.} Система за п. 1, яка відрізняється тим, що додатково містить модуль класичного нейромережевого аналізу (195), з'єднаний зі сховищем (200) та виконаний з можливістю оцінювання стану flow за розширеним вектором ознак \(\mathbf{f}_t \in \mathbb{R}^{22}\), що містить спектральні потужності у 5 діапазонах по 4 електродах, індекс складності задачі та індикатор якості сигналу, і повертає ймовірність стану flow \(p_{\text{flow}}^{\text{NN}}(t) \in [0,\,1]\) та бінарну мітку \(\text{FlowLabel}\), причому мітки зберігаються у сховищі для подальшого навчання квантової моделі.



\subsection{12. Гібридний квантово-класичний інференс}\label{ux433ux456ux431ux440ux438ux434ux43dux438ux439-ux43aux432ux430ux43dux442ux43eux432ux43e-ux43aux43bux430ux441ux438ux447ux43dux438ux439-ux456ux43dux444ux435ux440ux435ux43dux441}



\textbf{12.} Система за п. 11, яка відрізняється тим, що підтримує гібридний режим інференсу, в якому підсумкова ймовірність стану flow обчислюється як зважена комбінація квантової та класичної оцінок:



\[p_{\text{flow}}^{\text{hybrid}}(t) = \mu \cdot p_{\text{flow}}(t) + (1 - \mu) \cdot p_{\text{flow}}^{\text{NN}}(t)\]



де \(\mu \in [0,\,1]\) --- конфігурований коефіцієнт довіри до квантової моделі.



\subsection{13. Асинхронне збереження та евристичне маркування}\label{ux430ux441ux438ux43dux445ux440ux43eux43dux43dux435-ux437ux431ux435ux440ux435ux436ux435ux43dux43dux44f-ux442ux430-ux435ux432ux440ux438ux441ux442ux438ux447ux43dux435-ux43cux430ux440ux43aux443ux432ux430ux43dux43dux44f}



\textbf{13.} Система за п. 1, яка відрізняється тим, що додатково містить модуль асинхронного збереження даних (205), виконаний з можливістю прийому потокових EEG-вікон через обмежену чергу ємністю \(N_{\text{buf}}\) елементів, пакетного запису у сховище та автоматичного евристичного маркування стану flow за фронтальним \(\alpha/\theta\)-відношенням:



\[r(t) = \frac{P_\alpha(\text{AF7}, t) + P_\alpha(\text{AF8}, t)}{P_\theta(\text{AF7}, t) + P_\theta(\text{AF8}, t)}, \quad \text{FlowLabel}_{\text{heur}}(t) = \mathbb{1}[r(t) \geq 1.0]\]



причому запис у сховище виконується асинхронно відносно основного потоку обробки та обчислення CME, що забезпечує неблокуючий характер збереження.



\subsection{14. Маркування даних класичною моделлю для навчання квантової}\label{ux43cux430ux440ux43aux443ux432ux430ux43dux43dux44f-ux434ux430ux43dux438ux445-ux43aux43bux430ux441ux438ux447ux43dux43eux44e-ux43cux43eux434ux435ux43bux43bux44e-ux434ux43bux44f-ux43dux430ux432ux447ux430ux43dux43dux44f-ux43aux432ux430ux43dux442ux43eux432ux43eux457}



\textbf{14.} Спосіб за п. 2, який відрізняється тим, що додатково включає крок маркування набору даних, на якому: (i) зібрані вікна EEG-ознак передаються до модуля класичного нейромережевого аналізу (195) для отримання міток \(\text{FlowLabel}_{\text{NN}}(t)\) та ймовірностей \(p_{\text{flow}}^{\text{NN}}(t)\);



\begin{enumerate}

\def\labelenumi{(\roman{enumi})}

\setcounter{enumi}{1}

\tightlist

\item

  отримані мітки зберігаються у сховищі (200) та використовуються для навчання або перенавчання квантової моделі шляхом метаевристичної оптимізації (190).

\end{enumerate}


\end{document}
